\makeabstract
    {
    Pushing The Limits of Mentoring to the FURTHESST Extent: Improving STEM Persistence and Retention
    }
    {
    Julia Aram\textsuperscript{1}, Mu Jia Jiang\textsuperscript{1}, Kate Follette\textsuperscript{1}, Vonna Hemmler\textsuperscript{1}, Isaac Rosenthal\textsuperscript{1}
    }
    {
    \textsuperscript{1} Department of Physics \& Astronomy, Amherst College, Amherst, MA 
    }
    {
    Effective mentoring is one of the most important factors in improving students{\textquoteright} self-efficacy and perception of their skills, their science identity and interest in STEM, and their sense of belonging and integration within their science community, which can contribute to their persistence in the STEM field. Despite a growing body of literature affirming the benefits of mentoring, there still needs to be further analysis conducted on why and how mentoring is such an impactful practice, especially at the undergraduate level. In the Facilitating Undergraduate Research Through Holistic and Extensive Soft Skills Training (FURTHESST) Project, we explore the effects of mentoring relationships and current lab practices in helping students navigate this {\textquotedblleft}hidden curriculum{\textquotedblright} of scientific research.
    }
    {
    This summer, our team explored how detailed mentoring and lab practices have impacted long-term alumni outcomes. Through reading and annotating relevant literature, quantitatively analyzing data from alumni surveys, and qualitatively coding open-ended survey responses, we examined the lasting effects that key mentoring and lab systems have had on lab alumni since 2017. 
    }
    {
    In our quantitative analysis, we found that lab alumni who have spent longer time in the lab (8 or more time periods) were more likely to report that they Strongly Agreed or Agreed to feeling skill growth or gain in certain hard skills. They were also more likely to report that they Strongly Agreed or Agreed to feeling like a member of the lab community at Amherst and beyond. In our qualitative analysis, key themes included direct mentor-mentee interactions, interconnected outcomes, peer mentoring, and psychosocial and instrumental support.
    }
    {
    This project is critical for understanding the processes, emerging themes, and key findings within the Follette Lab{\textquoteright}s mentoring and practices that we can continue to develop, replicate, and implement to assist students in succeeding in open-ended research. 
    }

    \newpage
    